\section{原题H：用简短题目检查概念是否建立}
本章保留原稿H类题目，并把选项背后的理由展开。选择题的答案往往依赖题目的简化条件；理解这些条件，比只记选项字母更可靠。

\subsection{H1：椒盐噪声为什么优先使用中值滤波}
\label{original:H1}
\textbf{原题：}下列哪个模板可用于去除椒盐噪声？A. $3\times3$均值；B.中值滤波；C.拉普拉斯；D.高提升。

在本题的典型孤立脉冲噪声情形下，选B。椒盐噪声表现为个别异常亮点或暗点。中值滤波先把邻域的数按大小排序，再取中间的数；只要异常值没有占据足够多的位置，最大的255或最小的0就不会被选为中间值。例如$[98,99,100,101,255]$的中值是100，而均值为130.6。均值会受到极大值的拉动。

拉普拉斯和高提升主要突出变化，可能进一步放大噪声。中值也不是对任意密度的椒盐噪声都有效：异常点太多或目标太细时，仍可能失效或删除真实细节。

\subsection{H2：点处理与邻域处理的分界}
\label{original:H2}
\textbf{原题：}下列属于空间域点处理的是？A.均值滤波；B.中值滤波；C.伽马变换；D.拉普拉斯锐化。

选C。判断只需看“计算某个输出像素时，需要读哪些输入”。伽马变换先把该像素灰度归一化为$u\in[0,1]$，再计算$v=u^\gamma$并缩放回目标灰度范围；这一过程只读取该像素本身的数值，因此属于点处理。均值要读取周围多个像素，中值要对邻域排序，拉普拉斯要比较中心与邻居，都属于邻域处理。

\subsection{H3：二阶差分的双响应}
\label{original:H3}
\textbf{原题：}对灰度阶跃最能产生“双响应”的是？A.一阶微分；B.二阶微分；C.均值滤波；D.中值滤波。

选B。考虑从0突然变到10的序列$[0,0,0,10,10,10]$。一阶差分比较相邻两项，在跳变处产生一个大的变化量。再对这一阶变化作差分，一侧由0升到较大值，另一侧由较大值降回0，于是得到一正一负的一对响应。二阶差分强调的是“变化量本身怎样变化”，而不只是灰度值是否高。

\subsection{H4：保边平滑的条件}
\label{original:H4}
\textbf{原题：}能同时“平滑噪声且较少模糊边缘”的是？A.均值；B.高斯低通；C.中值；D.理想低通。

按这组入门选项和脉冲噪声背景，选C。边缘两侧是不同的真实灰度，均值会把两侧混合为中间灰度，使边界变缓；中值选取邻域已有灰度，往往能在去掉孤立异常点的同时保留阶跃。不过，这并不是所有噪声类型下的统一排名。中值会影响细线和小目标；高斯滤波也有自己的适用场景。

\subsection{H5：8比特为什么是256级}
\label{original:H5}
\textbf{原题：}$8\,\mathrm{bit}$灰度图像最多能表示多少级灰度？

每个比特只有0和1两种状态。两个比特有$2\times2=4$种组合，八个比特有$2^8=256$种组合。因此最多表示256级灰度，常用整数编码为0至255。最大编码是255，而级数是256，因为0也占一个级别。

\subsection{H6：均衡化为什么使用累计分布}
\label{original:H6}
\textbf{原题：}直方图均衡的变换函数是什么？

基本构造是累积分布函数（Cumulative Distribution Function, CDF），再乘目标最大灰度并按规定取整。某个灰度的CDF表示“灰度不超过它的像素占了多大比例”。因此它把输入灰度的位置，改写为其在整幅图中的累计比例位置。连续理想模型中，这能把灰度概率重新分配；离散图像中多个输入等级可能合并，输出并不保证每个灰度都有相同像素数。

\subsection{H7：相同灰度为何看起来不同}
\label{original:H7}
\textbf{原题：}“同时对比”说明感知亮度依赖于什么？

依赖周围的背景亮度。同一块灰色放在黑背景中常显得较亮，放在白背景中常显得较暗。两块中心区域的数字可以完全相同，变化来自视觉系统的相对判断。因此图像中的一个灰度数值和人眼感到的亮度不能简单画等号。

\subsection{H8：混合邻接限制哪一种对角连接}
\label{original:H8}
\textbf{原题：}$m$邻接主要用于解决什么问题？

混合邻接（Mixed Adjacency）用更严格的条件处理对角像素之间的连接，避免8邻接引入的某些局部多路连接歧义。两个候选像素若是对角邻居，只有它们的共同4邻域中没有属于指定前景值集合的像素时，才允许这条对角连接。

这里检查的是\textbf{共同4邻域}，不是搜索整幅图中是否存在任意一条4连通路径；也不能把规则概括为“整张图的路径从此唯一”。

\subsection{H9：系数和为零会去掉什么}
\label{original:H9}
\textbf{原题：}纯锐化（高通）掩模的系数和是多少？

对题目所指的去直流导数型高通核，系数和为0。若邻域全是常数$c$，输出为$c\sum w$；当$\sum w=0$时，均匀背景输出0，只对变化产生响应。

但“最终锐化图像”通常还需要加回原图。加回之后的有效核可以具有系数和1，以保留均匀区域的灰度。必须区分\textbf{提取细节的高通核}与\textbf{输出增强图像的完整锐化核}。

\subsection{H10：CLAHE限制的是哪一个步骤}
\label{original:H10}
\textbf{原题：}CLAHE相比AHE的关键改进是什么？

限制局部直方图的峰高，即直方图裁剪（Histogram Clipping），并在局部映射之间进行插值以减轻块边界。插值也可以用于AHE，不能把它当成CLAHE独有的区别；名称中的“限制对比度”主要对应裁剪机制。过高的峰会使局部累计映射变化过快，从而把很小的灰度差异放得过大。限制峰高可以减轻噪声过度增强，但不能保证噪声完全不被放大。

\subsection{H11：频谱中心化是重排坐标}
\label{original:H11}
\textbf{原题：}频域滤波前把图像乘$(-1)^{x+y}$有什么作用？

在常见的偶数尺寸离散图像模型中，它把频谱的直流分量移到中心，使低通和高通滤波器便于按照“到中心的距离”描述。这里$(-1)^{x+y}$按棋盘格在$+1$与$-1$之间交替。它不是消除噪声，也不是改变截止频率本身。奇数尺寸与具体软件索引需要额外留意；不能把这一乘法与任意尺寸的\texttt{fftshift}无条件等同。

\subsection{H12：突变截止为何容易振铃}
\label{original:H12}
\textbf{原题：}下列最易产生振铃的是？A.高斯低通；B.巴特沃斯低通（$n=2$）；C.理想低通；D.均值滤波。

在本题的典型滤波器比较中，选C。理想低通在截止处从1突然跳到0，对应的空间响应有交替的旁瓣。旁瓣与图像阶跃边缘叠加，容易产生边缘附近的明暗振荡。高斯低通平滑衰减，其理想空间核没有这种正负旁瓣；巴特沃斯的阶数越高，通常越接近突变截止，振铃风险也随之增加。

\subsection{H13：周期条纹与成对频谱峰}
\label{original:H13}
\textbf{原题：}去除图像上周期性条纹干扰宜用什么方法？

当干扰集中在可识别的频率位置时，可以使用频域陷波滤波（Notch Filtering），或适当的带阻滤波（Band-reject Filtering）。实值图像的周期性正弦条纹通常对应关于频谱中心成对的峰，因此陷波位置应成对考虑。

选择抑制区域时必须兼顾真实图像细节。若真实栅栏、织物等结构恰好占据相近频率，滤波也可能把它们削弱。频谱中的亮点不是自动等于噪声，必须结合原图判断。

\section{原稿与详解版之间的对应及更正}
本书的内容范围来自提供的黑白复习讲义，原题编号用于对应，不代表所有题目都已由教师确认为考试原题。原稿未提供的题图不会被伪造为原图；相应图形仅作为明确标注的教学示意。

主要更正包括：视场与物方像素大小的数值一致性；拉普拉斯公式和模板的正负号；A2中二阶差分核的性质；中值滤波的非线性与保边条件；高提升滤波的亮度变化；离散直方图均衡与BBHE的适用条件；DFT线性卷积所需的补零尺寸；振幅0.5对应约$-6.02$ dB而非$-3$ dB；同态滤波的常数增益在原图域通常对应幂变换。

这些更正已经融入相应知识解释和解题过程，不再要求同时记忆相互冲突的两套答案。对于工程案例，保留题目提供的假设，但把“题设条件”“理想计算”和“实际设备是否满足”分开说明。
